% !TEX root = ../main.tex
\documentclass[12pt,a4paper,twoside,openany]{book}
% Version écran : chapitres sur n'importe quelle page, pas de page blanche.
% Pour la version IMPRIMÉE reliée : remettre "openright", supprimer la ligne
% ci-dessous, et remettre \cleardoublepage à la place des \clearpage du
% début de main.tex. La numérotation reste en chiffres arabes dans les deux cas.
\let\cleardoublepage\clearpage

% --- Langue et encodage ---
\usepackage[utf8]{inputenc}
\usepackage[T1]{fontenc}
\usepackage{lmodern}
\usepackage[french]{babel}
\usepackage{csquotes}

% --- Mise en page ---
\usepackage[top=2.5cm,bottom=2.5cm,inner=3cm,outer=2.5cm]{geometry}
\usepackage{setspace}
\onehalfspacing
\usepackage{microtype}
\usepackage{fancyhdr}
\usepackage[explicit]{titlesec}
\usepackage{titletoc}
\usepackage{emptypage}

% --- Couleurs (charte Al Wassit : anthracite et jaune) ---
\usepackage[table,dvipsnames]{xcolor}
\definecolor{anthracite}{HTML}{111827}
\definecolor{jauneaw}{HTML}{F5B700}
\definecolor{grisclair}{HTML}{F3F4F6}
\definecolor{griscode}{HTML}{F8F8F8}
\definecolor{rougeatodo}{HTML}{C00000}

% --- Figures, tableaux, listes ---
\usepackage{graphicx}
\graphicspath{{figures/}}
\usepackage{float}
\usepackage[font=small,labelfont=bf,labelsep=period]{caption}
\usepackage{subcaption}
\usepackage{booktabs}
\usepackage{tabularx}
\usepackage{longtable}
\usepackage{array}
\usepackage{multirow}
\usepackage{enumitem}
\setlist{itemsep=2pt,topsep=4pt}
\usepackage{tikz}
\usetikzlibrary{arrows.meta,positioning,shapes.geometric,fit,calc}
\usepackage{pdflscape}
\usepackage{adjustbox}

% --- Mathématiques / symboles ---
\usepackage{amsmath,amssymb}
\usepackage{pifont}

% --- Code source ---
\usepackage{listings}
\lstdefinelanguage{JSONx}{
  morestring=[b]",
  literate={:}{{{\color{anthracite}:}}}1,
  sensitive=false
}
\lstset{
  basicstyle=\ttfamily\footnotesize,
  backgroundcolor=\color{griscode},
  frame=single,
  rulecolor=\color{gray!50},
  breaklines=true,
  breakatwhitespace=false,
  columns=fullflexible,
  keepspaces=true,
  showstringspaces=false,
  numbers=left,
  numberstyle=\tiny\color{gray},
  numbersep=6pt,
  xleftmargin=1.2em,
  framexleftmargin=1em,
  commentstyle=\color{gray}\itshape,
  keywordstyle=\color{anthracite}\bfseries,
  stringstyle=\color{brown!80!black},
  captionpos=b,
  literate=
    {é}{{\'e}}1 {è}{{\`e}}1 {ê}{{\^e}}1 {à}{{\`a}}1 {â}{{\^a}}1
    {ç}{{\c c}}1 {î}{{\^i}}1 {ô}{{\^o}}1 {û}{{\^u}}1 {ù}{{\`u}}1
    {É}{{\'E}}1 {À}{{\`A}}1 {ï}{{\"i}}1
}
\lstdefinelanguage{JavaScriptx}{
  keywords={const,let,var,function,return,if,else,async,await,import,export,from,new,throw,try,catch,default,true,false,null,for,of},
  sensitive=true,
  morecomment=[l]{//},
  morecomment=[s]{/*}{*/},
  morestring=[b]',
  morestring=[b]",
  morestring=[b]`
}
\lstdefinelanguage{TypeScriptx}{
  language=JavaScriptx,
  morekeywords={interface,type,class,implements,readonly,private,public,inject,Injectable}
}

% --- Bibliographie ---
\lstset{literate={é}{{\'e}}1 {è}{{\`e}}1 {à}{{\`a}}1 {ç}{{\c c}}1 {ê}{{\^e}}1 {É}{{\'E}}1}

\usepackage[backend=biber,style=numeric,sorting=none,maxbibnames=4]{biblatex}
\addbibresource{bibliographie.bib}

% --- Liens (en dernier ou presque) ---
\usepackage[hidelinks,colorlinks=false,pdfusetitle,pdfencoding=auto,unicode]{hyperref}
\usepackage[nameinlink,noabbrev]{cleveref}
\usepackage{url}
\urlstyle{same}

% --- Glossaire / sigles ---
\usepackage[acronym,nomain,toc,nonumberlist,style=long3col]{glossaries-extra}

% --- Annotations de rédaction (à désactiver pour la version finale) ---
% Mettre \afficheratodotrue en \afficheratodofalse pour masquer les repères rouges.
\newif\ifafficheratodo
\afficheratodotrue
\newcommand{\todo}[1]{%
  \ifafficheratodo\textcolor{rougeatodo}{\textbf{[À COMPLÉTER : #1]}}\fi}

% Emplacement réservé pour une capture ou un schéma encore manquant.
% Usage : \figmanquante{label}{Légende}{Description précise de ce qu'il faut capturer}
\newcommand{\figmanquante}[3]{%
  \begin{figure}[htbp]
    \centering
    \fbox{\parbox{0.88\textwidth}{\centering\vspace{1.2cm}%
      \textcolor{rougeatodo}{\textbf{FIGURE À INSÉRER}}\\[4pt]
      \small #3\vspace{1.2cm}}}
    \caption{#2}\label{#1}
  \end{figure}}

% Figure standard (largeur en fraction de \textwidth)
\newcommand{\fig}[4][0.9]{%
  \begin{figure}[htbp]
    \centering
    \includegraphics[width=#1\textwidth]{#2}
    \caption{#3}\label{#4}
  \end{figure}}

% Source d'une figure ou d'un tableau
\newcommand{\source}[1]{\par\vspace{-2pt}{\footnotesize\textit{Source : #1}}\par}

% Raccourcis typographiques
\newcommand{\nom}[1]{\textsc{#1}}
\newcommand{\code}[1]{\texttt{\small #1}}
\newcommand{\route}[1]{{\small\path{#1}}}
\newcommand{\xapp}{Xapp}

% --- Titres de chapitres ---
\titleformat{\chapter}[display]
  {\normalfont\bfseries\color{anthracite}}
  {\filleft\Large\textsc{\chaptertitlename}~\Huge\thechapter}
  {10pt}
  {\titlerule[2pt]\vspace{8pt}\filleft\Huge #1}
\titlespacing*{\chapter}{0pt}{0pt}{30pt}
\titleformat{name=\chapter,numberless}[display]
  {\normalfont\bfseries\color{anthracite}}
  {}
  {0pt}
  {\titlerule[2pt]\vspace{8pt}\filleft\Huge #1}
\titleformat{\section}{\normalfont\Large\bfseries\color{anthracite}}{\thesection}{0.8em}{#1}
\titleformat{\subsection}{\normalfont\large\bfseries\color{anthracite}}{\thesubsection}{0.8em}{#1}
\titleformat{\subsubsection}{\normalfont\normalsize\bfseries\color{anthracite}}{\thesubsubsection}{0.8em}{#1}
\setcounter{secnumdepth}{3}
\setcounter{tocdepth}{2}

% --- En-têtes et pieds de page ---
\pagestyle{fancy}
\fancyhf{}
\fancyhead[LE]{\small\nouppercase{\leftmark}}
\fancyhead[RO]{\small\nouppercase{\rightmark}}
\fancyfoot[LE,RO]{\thepage}
\fancyfoot[RE,LO]{\footnotesize\titrecourt}
\renewcommand{\headrulewidth}{0.4pt}
\renewcommand{\footrulewidth}{0.2pt}
\fancypagestyle{plain}{%
  \fancyhf{}\fancyfoot[LE,RO]{\thepage}\renewcommand{\headrulewidth}{0pt}}

% --- Environnement encadré pour les définitions / remarques ---
\usepackage[most]{tcolorbox}
\newtcolorbox{remarque}[1][Remarque]{
  colback=grisclair,colframe=jauneaw,fonttitle=\bfseries,title=#1,
  boxrule=0.6pt,left=6pt,right=6pt,top=4pt,bottom=4pt,breakable}

\setlength{\parindent}{1.2em}
\setlength{\parskip}{3pt}
\raggedbottom
\frenchbsetup{StandardItemLabels=true}
